%% Apêndices

\begin{apendicesenv}

\chapter{Dicionário de Atributos do Conjunto de Dados Sintéticos}
\label{ap:dicionario}

A Tabela~\ref{tab:dicionario_atributos} apresenta o dicionário completo dos atributos gerados pelo gerador de dados sintéticos, referenciado no capítulo de Metodologia.

\begin{table}[htb]
\centering
\caption{Dicionário descritivo dos atributos do conjunto de dados sintéticos.}
\label{tab:dicionario_atributos}
\small
\begin{tabular}{llp{6.4cm}}
\toprule
\textbf{Atributo} & \textbf{Tipo} & \textbf{Descrição} \\
\midrule
\texttt{serie} & Numérico (1--3) & Série do ensino médio em curso. \\
\texttt{idade} & Numérico (14--24) & Idade do estudante, com assimetria positiva decorrente da distorção idade-série. \\
\texttt{genero} & Categórico (binário) & Gênero autodeclarado (paridade $\sim$50/50). \\
\texttt{cor\_raca} & Categórico (nominal) & Cor/raça autodeclarada (Parda, Preta, Branca, Outra). \\
\texttt{estado\_civil} & Categórico (nominal) & Estado civil (Solteiro, Casado/União, Outro). \\
\texttt{renda\_per\_capita} & Categórico (nominal) & Faixa de renda familiar \textit{per capita} (até 0,5 SM; 0,5--1 SM; acima de 1 SM). \\
\texttt{trabalha} & Categórico (binário) & Indica se o estudante concilia trabalho e estudo. \\
\texttt{acesso\_internet} & Categórico (binário) & Acesso doméstico à internet. \\
\texttt{possui\_computador} & Categórico (binário) & Disponibilidade de computador para os estudos (proxy da exclusão digital qualitativa). \\
\texttt{tipo\_escola\_origem} & Categórico (binário) & Rede de origem no ensino fundamental (pública ou privada). \\
\texttt{distorcao\_idade\_serie} & Categórico (binário) & Indica atraso escolar de dois ou mais anos. \\
\texttt{assistencia\_pe\_de\_meia} & Categórico (binário) & Beneficiário efetivo do programa Pé-de-Meia (derivado por regra de elegibilidade). \\
\texttt{nota\_media} & Numérico (0--10) & Média das notas parciais do 1º semestre letivo. \\
\texttt{frequencia} & Numérico (\%) & Percentual de frequência às aulas acumulado no 1º semestre. \\
\texttt{media\_atividades} & Numérico (0--10) & Média das notas obtidas nas atividades avaliativas do 1º semestre, normalizada para a escala 0--10 (ex.: atividade valendo 2,0 pontos com nota 1,0 equivale a 5,0). \\
\texttt{taxa\_entrega\_atividades} & Numérico (\%) & Proporção de atividades entregues no 1º semestre (sinal de regularidade do engajamento). \\
\texttt{reprovacoes} & Numérico (0--3) & Reprovações acumuladas no ensino médio. \\
\texttt{situacao} & Categórico (alvo) & Tendência de desempenho projetada para o final do ano letivo, rotulada pelos critérios da LDB: Aprovado, Em Risco ou Reprovado. \\
\bottomrule
\end{tabular}
\source{Elaborado pelo autor (2026).}
\end{table}

\end{apendicesenv}
